% Part: proof-theory
% Chapter: normalization
% Section: translations

\documentclass[../../../include/open-logic-section]{subfiles}

\begin{document}

\olfileid{pt}{nor}{tr}

\olsection{正规 \Log{N2i}-推导与 \Log{G2i} 之间的互译}

\Log{N2i} 中的推导可以翻译到 \Log{G2i} 中的推导，反之亦然。首先，将无切割的
\Log{G2i} 推导翻译过来，得到的是正规的 \Log{N2i} 推导。

为了陈述这个结果，我们需要一些术语，以便更容易地将~\Log{N2i} 和~\Log{G2i} 中相继式的前件联系起来：如果对于每个公式~$!A$，$\Gamma'$ 中形如 $x : !A$ 的元素个数小于或等于~$!A$ 在~$\Gamma$ 中的多重度，则称带标签公式的集合~$\Gamma'$ \emph{对应}于不带标签公式的多重集~$\Gamma$。

\begin{cor}
如果 $\Log{G2i} \Proves \Gamma \Sequent \Delta$，那么存在 $\Gamma' \Sequent !C$ 的正规
\Log{N2c}-推导~$\delta$，其中带标签公式的集合~$\Gamma'$ 对应于多重集~$\Gamma$；也就是说，对于每个公式~$!A$，$\Gamma'$ 中形如 $x : !A$ 的元素个数小于或等于~$!A$ 在~$\Gamma$ 中的多重度（即~$!A$ 的拷贝数）。
\end{cor}

\begin{proof}
检查 \olref[nat][tgi]{prop:G2i-to-N2i} 的证明可知：我们将引入规则加到推导的末尾，而消去规则仅出现在推导的顶部，因此绝不会出现引入规则在消去规则之上的情况。
\end{proof}

然而，\olref[nat][tgi]{prop:G2i-to-N2i} 中给出的 \Cut{} 规则的翻译，可能会在
\Log{N2i} 推导~$\delta$ 中产生切割。如果一个 \Log{G2i} 推导以~\Cut{} 结束，且其直接子推导分别为 $\pi_1$ 和 $\pi_2$，分别推导出 $\Gamma_1
\Sequent !A$ 和 $!A, \Gamma_2 \Sequent \Delta_2$，那么我们将 $\pi_1$ 的翻译 $\delta_1$ 嫁接到 $\pi_2$ 的翻译 $\delta_2$ 中出现的公理 $x:!A
\Sequent !A$ 上。如果 $\delta_1$ 以一个主公式为~$!A$ 的引入规则结束，并且 $\delta_2$ 中的公理 $x:!A \Sequent !A$ 是某个消去规则的主前提，这就会导致 $\delta$ 中出现一个切割。

也可以将 \Log{N2i} 中的推导翻译到~\Log{G2i} 中的推导。在 \olref[nat][tni]{prop:N2-to-G2} 的证明中，\Elim{\land}、\Elim{\lif}、\Elim{\lnot}、
\Elim{\lforall} 的翻译会在得到的 \Log{G2i} 推导中产生~$\Cut$ 推理。然而，如果我们从一个正规的 \Log{N2i} 推导出发，这是可以避免的。

如果 $S_1$ 是公理，并且只要 $S_i$ 是某个推理的主前提，就有 $S_{i+1}$ 是该推理的结论，我们就称 \Log{N2i} 推导~$\delta$ 中相继式出现的一个列表 $S_1$, \dots,~$S_n$ 为一个\emph{分支}。换句话说，一个分支在 $\delta$ 中从公理开始向下追踪相继式，但在遇到消去规则的次前提时停止。$\delta$ 的一个\emph{主分支}是满足 $S_n$ 为终结相继式的分支。每个推导~$\delta$ 必须至少有一个主分支，因为每条规则都至少有一个主前提（只有 \Intro{\land} 有两个）。

\begin{prop}\ollabel{prop:normal-N2i-to-G2i} 如果 $\delta$ 是 $\Gamma \Sequent !A$ 的正规 $\Log{N2c}$-或 $\Log{N2i}$-推导，那么存在 $\Gamma' \Sequent \Delta$ 的（无切割的）$\Log{G2c}$-推导（或 $\Log{G2i}$-推导）~$\pi$，其中 $\Gamma'$ 是从 $\Gamma$ 去掉标签后得到的公式多重集，并且如果 $!A$ 不是~$\lfalse$，则 $\Delta = \{!A\}$，若是则 $\Delta = \emptyset$。
\end{prop}

\begin{proof}
这一次，我们对 $\Gamma \Sequent !A$ 的推导~$\delta$ 中的推理数目进行归纳。
如果~$\delta$ 中没有推理，那么它是一个公理 $x:!A \Sequent !A$，
而 $\pi$ 就是相应的公理 $!A \Sequent !A$，
或者当 $!A \ident \lfalse$ 时为 $\lfalse \Sequent \ $。

如果 $\delta$ 以一条推理规则结束，我们区分情况：
\begin{enumerate}
  \item 如果 $R$ 是一条引入规则或~\FalseInt，
  翻译过程与 \olref[nat][tni]{prop:N2-to-G2} 的证明完全相同，且不需要~\Cut{}。

  如果 $\delta$ 的最后一条规则~$R$ 是一条消去规则，则观察以下事实：
  \begin{enumerate}
    \item $\delta$ 的主分支不经过任何引入规则，否则主分支上就会有一个引入规则或~\FalseInt{} 之后紧跟一个消去规则，从而 $\delta$ 不是正规的。
    \item 由于 $\delta$ 的主分支上没有引入规则，因此主分支只有一条。（只有 \Intro{\land} 有两个主前提，所以多个主分支必然经过一个 \Intro{\land} 规则的前提。）
    \item 如果主分支包含了 \Elim{\lor} 或 \Elim{\lexists} 的主前提，那么这样的规则只有一条，且它是最后一条推理，即~$R$。（否则，主分支上会包含一个 \Elim{\lor} 或 \Elim{\lexists} 规则，而该规则的结论又是某个消去规则的主前提，这样我们就可以应用一个置换转换，推导就不再是正规的。）
    \item 除了 \Elim{\lor} 和 \Elim{\lexists} 之外，没有消去规则会解除假设，即改变相继式的左上下文。\Elim{\lor} 和 \Elim{\lexists} 只解除其次前提之上的假设，即只改变其次前提的左上下文。换言之，如果 $\delta$ 主分支上的相继式 $S_i$ 具有上下文~$\Gamma_1$，那么以 $S_i$ 为主前提的推理的结论 $S_{i+1}$ 具有上下文 $\Gamma_1' \supseteq \Gamma_1$。
    \item 因此，如果 $x: !C \Sequent !C$ 是主分支中最顶端的相继式~$S_1$，则 $x:!C \in \Gamma$。
  \end{enumerate}
  我们现在根据~$!C$ 的形式来区分情况。由于主分支上的每一个 $S_i$ 都是某个消去规则的主前提，这对于 $x:!C \Sequent !C$ 同样成立。

  \item $!C \ident !D \land !E$。那么 $\delta$ 的形式为
  \[
  \Axiom$x:!D \land !E \fCenter !D \land !E$
  \RightLabel{\Elim{\land}[1]}
  \UnaryInf$x:!D \land !E \fCenter !D$
  \Deduce$x:!D \land !E, \Gamma_1 \fCenter !A$
  \DisplayProof
  \]
  由于包含 $x:!D \land !E \Sequent !D$ 的主分支不含 \Intro{\lif} 或 \Intro{\lnot} 的主前提，也不含 \Elim{\lor} 或 \Elim{\lexists} 的次前提，$x:!D \land !E$ 中的上下文公式在主分支上保持不变。这解释了为何终结相继式的上下文必须包含~$x:!D \land !E$。
  此外，我们可以通过将结束于 \Elim{\land} 的子推导替换为 $x:!D \Sequent !D$，
  并将主分支上每一个相继式上下文中的 $x:!D \land !E$ 替换为 $x:!D$，从而将推导 $\delta$ 转化为推导~$\delta_1$，即把
  \[
  \bottomAlignProof
  \Axiom$x:!D \land !E \fCenter !D \land !E$
  \RightLabel{\Elim{\land}[1]}
  \UnaryInf$x:!D \land !E \fCenter !D$
  \Deduce$x:!D \land !E, \Gamma_1 \fCenter !A$
  \DisplayProof
  \quad\text{ 转化为 }\quad
  \bottomAlignProof
  \Axiom$x:!D \fCenter !D$
  \Deduce$x:!D, \Gamma_1 \fCenter !A$
  \DisplayProof
  \]
  归纳假设可应用于 $\delta_1$，因为 $\delta$ 中的推理数目等于 $\delta_1$ 中的推理数目加~$1$。

  由归纳假设，存在 \Log{G2i}-推导 $\pi_1$，其结论为
  $x:!D, \Gamma_1 \Sequent !A$。
  我们通过应用 \LeftR{\land} 来得到推导~$\pi$：
  \[
  \AxiomC{}
  \RightLabel{$\pi_1$}
  \Deduce$!D, \Gamma_1' \fCenter !A$
  \RightLabel{\LeftR{\land}}
  \UnaryInf$!D \land !E, \Gamma_1' \fCenter !A$
  \DisplayProof
  \]
  如果 $!A \ident \lfalse$，我们插入一个 \RightR{\Weakening}，例如：
  \[
  \AxiomC{}
  \RightLabel{$\pi_1$}
  \Deduce$!D, \Gamma_1' \fCenter $
  \RightLabel{\LeftR{\land}}
  \UnaryInf$!D \land !E, \Gamma_1' \fCenter $
  \RightLabel{\RightR{\Weakening}}
  \UnaryInf$!D \land !E, \Gamma_1' \fCenter !A$
  \DisplayProof
  \]
  \item 如果 $!C \ident !D \lor !E$，
  那么它必须是 \Elim{\lor} 的主前提，并且此推理必须是主分支上的最后一条推理~$R$。
  换言之，$\delta$ 的形式为：
  \[
  \Axiom$x:!D \lor !E \fCenter !D \lor !E$
  \AxiomC{}
  \RightLabel{$\delta_1$}
  \Deduce$y:!D, \Gamma_1 \fCenter!A$
  \AxiomC{}
  \RightLabel{$\delta_2$}
  \Deduce$z:!E, \Gamma_2 \fCenter !A$
  \RightLabel{\Elim{\lor}}
  \TrinaryInf$x:!D \lor !E, \Gamma_1, \Gamma_2 \fCenter !A$
  \DisplayProof
  \]
  归纳假设可应用于推导 $\delta_1$ 和 $\delta_2$；
  我们得到 \Log{G2i}-推导 $\pi_1$ 和 $\pi_2$，其结论分别为 $!D, \Gamma_1' \Sequent !A$ 和 $!E, \Gamma_2' \Sequent !A$。
  我们应用 $\LeftR{\lor}$ 来得到~$\pi$：
  \[
  \AxiomC{}
  \RightLabel{$\pi_1$}
  \Deduce$!D, \Gamma_1' \fCenter !A$
  \AxiomC{}
  \RightLabel{$\pi_2$}
  \Deduce$!E, \Gamma_2' \fCenter !A$
  \RightLabel{\LeftR{\lor}}
  \BinaryInf$!D \lor !E, \Gamma_1', \Gamma_2' \fCenter !A$
  \DisplayProof
  \]
  如果 \Elim{\lor} 没有解除 $y: !D$ 或 $z: !E$（即它们未出现在次前提的上下文中），
  或者若~$!A$ 为 $\lfalse$，我们就需要添加一些 \LeftR{\Weakening} 和 \RightR{\Weakening}。
  例如：
  \[
  \AxiomC{}
  \RightLabel{$\pi_1$}
  \Deduce$\Gamma_1' \fCenter $
  \RightLabel{\LeftR{\Weakening}}
  \UnaryInf$!D, \Gamma_1' \fCenter $
  \AxiomC{}
  \RightLabel{$\pi_2$}
  \Deduce$\Gamma_2' \fCenter $
  \RightLabel{\LeftR{\Weakening}}
  \UnaryInf$!E, \Gamma_2' \fCenter $
  \RightLabel{\LeftR{\lor}}
  \BinaryInf$!D \lor !E, \Gamma_1', \Gamma_2' \fCenter $
  \RightLabel{\RightR{\Weakening}}
  \UnaryInf$!D \lor !E, \Gamma_1', \Gamma_2' \fCenter !A$
  \DisplayProof
  \]

  \item $!C \ident !D \lif !E$。那么 $\delta$ 的形式为
  \[
  \Axiom$x:!D \lif !E \fCenter !D \lif !E$
  \AxiomC{}
  \RightLabel{$\delta_1$}
  \Deduce$\Gamma_1 \fCenter!D$
  \RightLabel{\Elim{\lif}}
  \BinaryInf$x:!D \lif !E, \Gamma_1 \fCenter !E$
  \RightLabel{$\delta_2$}
  \Deduce$x:!D \lif !E, \Gamma_1, \Gamma_2 \fCenter !A$
  \DisplayProof
  \]
  注意，由于包含 $x:!D \lif !E,
  \Gamma_1 \Sequent !E$ 的主分支不含 \Intro{\lif} 或 \Intro{\lnot} 的主前提，也不含 \Elim{\lor} 或 \Elim{\lexists} 的次前提，$x:!D \lif !E, \Gamma_1$ 中的上下文公式在主分支上保持不变。这解释了为何终结相继式的上下文必须包含~$x:!D \lif !E, \Gamma_1$。此外，我们可以通过将结束于 \Elim{\lif} 的子推导替换为 $x:!E \Sequent !E$，并将主分支上每个相继式上下文中的 $x:!D \lif !E,
  \Gamma_1$ 替换为 $x:!E$，从而将推导片段 $\delta_2$ 转化为一个正确的推导~$\delta_2'$，即把
  \[
  \bottomAlignProof
  \Axiom$x:!D \lif !E, \Gamma_1 \fCenter !E$
  \RightLabel{$\delta_2$}
  \Deduce$x:!D \lif !E, \Gamma_1, \Gamma_2 \fCenter !A$
  \DisplayProof
  \quad\text{ 转化为 }\quad
  \bottomAlignProof
  \Axiom$x:!E \fCenter !E$
  \RightLabel{$\delta_2'$}
  \Deduce$x:!E, \Gamma_2 \fCenter !A$
  \DisplayProof
  \]
  归纳假设可应用于 $\delta_1$ 和 $\delta_2'$，因为 $\delta$ 中的推理数目等于 $\delta_1$ 和~$\delta_2$ 中推理数目之和，再加上以计入主分支起始处 \Elim{\lif} 推理的那~$1$ 个推理。（如果该推理同时也是 $\delta$ 中的最后一条推理~$R$，那么 $\delta_1$ 包含 $0$ 个推理。）

  由归纳假设，存在 \Log{G2i}-推导 $\pi_1$（结论为 $\Gamma_1 \Sequent !D$）和 $\pi_2$（结论为 $!E, \Gamma_1' \Sequent !A$）。
  我们通过应用 \LeftR{\lif} 来得到推导~$\pi$：
  \[
  \AxiomC{}
  \RightLabel{$\pi_1$}
  \Deduce$\Gamma_1' \fCenter !D$
  \AxiomC{}
  \RightLabel{$\pi_2$}
  \Deduce$!E, \Gamma_2' \fCenter !A$
  \RightLabel{\LeftR{\lif}}
  \BinaryInf$!D \lif !E, \Gamma_1', \Gamma_2' \fCenter !A$
  \DisplayProof
  \]
  如果 $!D$ 和/或 $!A \ident \lfalse$，我们插入一个或两个 \RightR{\Weakening}，例如：
  \[
  \AxiomC{}
  \RightLabel{$\pi_1$}
  \Deduce$\Gamma_1' \fCenter $
  \RightLabel{\RightR{\Weakening}}
  \UnaryInf$\Gamma_1' \fCenter !D$
  \AxiomC{}
  \RightLabel{$\pi_2$}
  \Deduce$!E, \Gamma_2' \fCenter $
  \RightLabel{\RightR{\Weakening}}
  \UnaryInf$!E, \Gamma_2' \fCenter !A$
  \RightLabel{\LeftR{\lif}}
  \BinaryInf$!D \lif !E, \Gamma_1', \Gamma_2' \fCenter !A$
  \DisplayProof
  \]

\end{enumerate}
其余情况类似处理，留作练习。
\end{proof}

\begin{cor}
在任何正规的 \Log{N1i} 或 \Log{N2i} 推导中出现的每个公式，都是结束公式或结束相继式的子公式。
\end{cor}

\begin{proof}
根据 \olref{prop:normal-N2i-to-G2i}，每个正规的 \Log{N2i} 推导都可以被翻译成一个 \Log{G2i} 推导。检查该证明可知，在 \Log{N2i} 推导中出现的每个公式，也会出现在对应的 \Log{G2i} 推导中。\Log{G2i} 不含 \Cut{} 规则，因此具有子公式性质。
\end{proof}

\end{document}